\subsubsection{The actual polynomial highest-weight algebra on the same ray}
The positive rule in \fid{F34} can be realized inside its ordinary
highest-weight algebra. The following direct slice calculation is a
constructive supplement to that finding, without a priority claim for
the classical method.

\paragraph{Claim.}
For $k\ge2$ put $V=\Sym^k\C^2$ and choose its divided basis
\[
 v_i=x^{k-i}y^i/i!,\qquad 0\le i\le k.
\]
In the free symmetric algebra $A=\Sym(V)=\C[v_0,\ldots,v_k]$,
let $D(v_0)=0$, $D(v_i)=v_{i-1}$, and assign
$\deg v_i=1$, $\operatorname{ind}v_i=i$.
The ray algebra
\[
 B=\bigoplus_{N\ge0}
       \bigl(A_{\deg=N,\,\operatorname{ind}=N}\cap\ker D\bigr)
\]
is a polynomial ring on $k-1$ explicit generators $f_2,\ldots,f_k$
of degrees $2,\ldots,k$. Its actual degree-$N$ monomial basis is
indexed by $2a_2+\cdots+ka_k=N$, and
\[
 \dim B_N=[\Sym^N(\Sym^k\C^2):S_{((k-1)N,N)}\C^2]=c_k(N).
\]
The symbols $v_i$ are free algebra generators here; their displayed
binary forms specify a basis and action, rather than an identification
of $A$ with a subalgebra of $\C[x,y]$.

\begin{proof}
The raising operator $x\partial_y$ acts by $Dv_i=v_{i-1}$ in
the divided basis. A degree-$N$, index-$r$ monomial has torus
weight $(kN-r,r)$. In characteristic zero, the vectors of this
weight killed by the raising operator are its highest-weight
vectors. Since $k\ge2$, the weight $((k-1)N,N)$ is dominant,
and ordinary highest-weight decomposition proves the multiplicity
identity. This also agrees with the classical two-row specialization
in \cite[Lemma~1.3, p.~2]{PPS} used in \fid{F34}.

In $A[v_0^{-1}]$ put $t=v_1/v_0$, so $Dt=1$, and define
\[
 f_i=v_0^{i-1}\sum_{j=0}^i\frac{(-t)^j}{j!}v_{i-j},
       \qquad 2\le i\le k.
\]
These elements belong to $A$. Indeed the terms with $j\le i-1$
are
\[
 \frac{(-1)^j}{j!}v_0^{i-1-j}v_1^jv_{i-j},
\]
and the last term is $(-1)^iv_1^i/i!$.
Each has degree and index $i$. Differentiating the sum, the terms
from $D((-t)^j/j!)$ cancel those from $D(v_{i-j})$ after shifting
the index; the remaining endpoint has $D(v_0)=0$.
Thus $Df_i=0$. For example,
\[
 f_2=v_0v_2-\tfrac12v_1^2,\qquad
 f_3=v_0^2v_3-v_0v_1v_2+\tfrac13v_1^3.
\]

The formula for $f_i$ is triangular, with leading term
$v_0^{i-1}v_i$ and remaining terms involving only earlier generators.
The leading coefficient is a unit in the localization. Solving
successively for $v_1=tv_0,v_2,\ldots,v_k$ therefore gives an
invertible change of free coordinates,
\[
 A[v_0^{-1}]=\C[v_0,v_0^{-1},t,f_2,\ldots,f_k].
\]
In these coordinates $D=\partial/\partial t$. Characteristic zero
then gives
\[
 A[v_0^{-1}]^D=\C[v_0,v_0^{-1},f_2,\ldots,f_k].
\]
This assertion concerns the explicitly computed localization; no
description of the entire unlocalized invariant ring is assumed.

Every monomial $v_0^a\prod_{i=2}^k f_i^{a_i}$ in this Laurent
polynomial ring has degree minus index equal to $a$.
Algebraic independence and homogeneity therefore imply that its
component with degree equal to index has $a=0$ and is exactly
$\C[f_2,\ldots,f_k]$. All the $f_i$ already belong to $A$, so
intersecting with $A$ loses nothing. Hence
\begin{equation}\label{eq:actualrayalgebra}
 B=\C[f_2,\ldots,f_k],\qquad \deg f_i=i.
\end{equation}
Its displayed monomials form the asserted actual highest-weight
basis. The construction is defined over $\Q$ and then base changed
to $\C$; all divided bases, factorials and the derivative-kernel
argument use characteristic zero.
\end{proof}

This gives multiplication constants equal to one in this particular
ordinary algebra. Its $k-1$ generators and their formulas have size
polynomial in the numerical $k$, not in $\log k$. The construction
therefore retains both the exact positive simplex and the compact
dimension obstruction in \fid{F34}. It does not identify this
monomial basis with VIII's nonstandard multiplication or
comultiplication basis, and asserts no theta or Weil realization.

% Separate paragraph for Part III.
\paragraph{The ray algebra in Grothendieck's language.}
The actual scheme $\Spec B$ is affine $(k-1)$-space with positive
grading $2,\ldots,k$; $\operatorname{Proj}B$ is the corresponding
weighted projective scheme. The degree-$N$ graded coordinate piece
is the highest-weight multiplicity space just proved, so its Hilbert
series is $\prod_{i=2}^k(1-z^i)^{-1}$. The exponent simplex has
dimension $k-2$. These are exact statements about the affine
coordinate ring and its grading, without an identification of
coherent Hilbert multiplicities with arbitrary constructible
cohomology. The ordinary affine and graded-scheme dictionary applies;
no arithmetic weight, designated nonstandard basis, or archival
claim is supplied by this coordinate computation alone.
